\section*{Erd\H{o}s Problems \#256--\#258: Round 2 continuation}

\section*{Erd\H{o}s Problem \#256}

\subsection*{1) ROUND-2 OBJECTIVE}
Path \textbf{(C)}: obstruction/correction. Round 1 already gave a complete disproof of the specific subquestion
\[
\exists c>0\text{ such that }\log f(n)\gg n^c.
\]
So the only legitimate round-2 task here is to adversarially check that disproof, lock in the quantifiers, and record the minimal corrected conclusion.

\subsection*{2) Round-1 FOUNDATION USED}
I use the following Round-1 facts without re-proving them.
\begin{enumerate}
\item Belov--Konyagin's upper bound: there exists an absolute constant $C>0$ such that
\[
\log f(n)\le C(\log n)^4
\]
for all sufficiently large $n$.
\item The elementary comparison
\[
(\log n)^4=o(n^c)\qquad (n\to\infty)
\]
for every fixed $c>0$.
\item The Round-1 contradiction argument combining (1) and (2).
\end{enumerate}

\subsection*{3) NEW INSIGHT / TOOL (ROUND-2)}
There is no new mathematical ingredient needed. The only new point is an adversarial consistency check against the current Erd\H{o}s Problems summary, which now explicitly records that the specific power-lower-bound question has a negative answer.

\subsection*{4) ATTACK PLAN (ROUND-2)}
The only possible gap after Round 1 would have been a quantifier slip: one must rule out \emph{every} fixed exponent $c>0$, not just a particular one. The plan is therefore to re-run the contradiction with the quantifiers written explicitly.

\subsection*{5) WORK (ROUND-2)}
Fix any $c>0$. Suppose, for contradiction, that there exist $K>0$ and $n_0$ such that
\[
\log f(n)\ge K n^c\qquad (n\ge n_0).
\]
Round 1 already established that there is an absolute constant $C>0$ such that
\[
\log f(n)\le C(\log n)^4
\]
for all sufficiently large $n$. Since $(\log n)^4=o(n^c)$, there exists $n_1$ such that
\[
C(\log n)^4<K n^c\qquad (n\ge n_1).
\]
Hence, for all sufficiently large $n$,
\[
\log f(n)\le C(\log n)^4<K n^c\le \log f(n),
\]
which is impossible. Therefore no fixed $c>0$ can satisfy $\log f(n)\gg n^c$.

Thus the specific subquestion in Problem \#256 is fully settled in the negative.

\subsection*{6) ADVERSARIAL VERIFICATION}
\begin{enumerate}
\item \textbf{Quantifiers.} The argument starts with an \emph{arbitrary} fixed $c>0$ and shows that this $c$ cannot work. Hence it rules out all $c>0$.
\item \textbf{Dependence of constants.} The Belov--Konyagin constant $C$ is absolute; the contradiction only needs that $C$ is independent of $n$.
\item \textbf{Logical target.} The original question only asked whether a power-type lower bound $\log f(n)\gg n^c$ could hold. The argument disproves exactly that claim, not the entire surrounding estimation problem for $f(n)$.
\end{enumerate}
No gap remains.

\subsection*{7) FINAL}
\textbf{FULL SOLUTION --- COUNTEREXAMPLE / DISPROOF}

The specific claim
\[
\exists c>0\text{ such that }\log f(n)\gg n^c
\]
is false. Belov--Konyagin give the upper bound $\log f(n)\ll (\log n)^4$, and $(\log n)^4=o(n^c)$ for every fixed $c>0$. Hence no power lower bound of the form $\log f(n)\gg n^c$ can hold.

\subsection*{8) COMPLETION ESTIMATE (MANDATORY)}
COMPLETION: 100\%

\subsection*{9) REFERENCES}
\begin{enumerate}
\item A. S. Belov and S. V. Konyagin, \emph{An estimate of the free term of a non-negative trigonometric polynomial with integer coefficients}, Izv. Math. \textbf{60} (1996), 1123--1182.
\item T. F. Bloom (ed.), \emph{Erd\H{o}s Problem \#256}, \texttt{erdosproblems.com/256}, accessed 2026-03-07.
\end{enumerate}

